\begin{exercise}{83}
Show that there exists a functor $f\colon\mathcal C\to\mathcal D$ between $\infty$-categories
which is conservative, but does not satisfy the RLP with respect to $\Delta^1\to J$.
Hint: Consider the map $J\to\mathcal S(|J|)$.

\begin{proof}
Proposition~1.1.36 gives the realization--singular adjunction.
Write its unit as $\eta\colon J\to\mathcal S(|J|)$.
Every edge of $J$ is an equivalence, and $\mathcal S(|J|)$
is a Kan complex by Lemma~1.1.47.
Thus $\eta$ is a conservative functor between $\infty$-categories.
Let $i\colon\Delta^1\to J$ be the edge $0\to1$.

Using the tuple description from Exercise~\ref{ex:60}, define
$g\colon|J|\to[0,1]$ on a simplex with barycentric coordinates by
\[
 g\bigl([(a_0,\ldots,a_n),(t_0,\ldots,t_n)]\bigr)
 =\sum_{i=0}^n a_it_i,
 \qquad a_i\in\{0,1\}.
\]
For $\alpha\colon[m]\to[n]$, the realization identifications
preserve this expression because
\[
 \sum_{i=0}^m a_{\alpha(i)}t_i
 =\sum_{j=0}^n a_j\sum_{\alpha(i)=j}t_i.
\]
It therefore descends to $|J|$ and is continuous by the
quotient topology, since its restriction to every simplex
is affine (Definition~1.1.35).
Identify $|\Delta^1|$ with $[0,1]$ and put
$r=|i|g\colon|J|\to|J|$.
Since $g|i|(t)=t$, we have $r|i|=|i|$.
However, $r\neq\mathrm{id}_{|J|}$:
its image lies in the realized edge $|i|$, whereas the
nondegenerate simplex $(0,1,0)$ gives an open $2$-cell
disjoint from the $1$-skeleton (Corollary~1.1.41).

Let $b\colon J\to\mathcal S(|J|)$ be the adjunct of $r$.
Naturality of the adjunction and $r|i|=|i|$ make the square
\[
\begin{tikzcd}
 \Delta^1 \arrow[r,"i"] \arrow[d,"i"']
   & J \arrow[d,"\eta"] \\
 J \arrow[r,"b"'] & \mathcal S(|J|)
\end{tikzcd}
\]
commute.
Any diagonal $d\colon J\to J$ must satisfy $di=i$,
so it fixes both vertices.
Every simplex of $J$ is determined by its tuple of vertices,
hence $d=\mathrm{id}_J$.
The other triangle would then give $b=\eta$.
Taking adjuncts would force $r=\mathrm{id}_{|J|}$,
a contradiction.
Thus this square has no lift.
\end{proof}
\end{exercise}
